\section{Introduction}\label{sec:introduction}
%% ----------------------------------------------------------------

LLM-based automated program repair has advanced rapidly, with benchmarks such as SWE-bench~\cite{swebench2024} and its human-validated \emph{Verified} subset becoming the standard yardstick for repository-level bug fixing. In this task an LLM receives a natural-language issue report and must produce a patch that passes the project's test suite---a setting that demands not only code-generation ability but also accurate \emph{localization} of the relevant source files and functions within repositories that may contain thousands of files. Context construction---deciding which code to show the LLM---is therefore a critical bottleneck, and the strategies adopted by leading systems such as SWE-Agent~\cite{yang2024sweagent}, AutoCodeRover~\cite{zhang2024autocoderover}, and Agentless~\cite{xia2024agentless} differ primarily in \emph{how} they select context rather than in the underlying language model.

The prevailing approach to context construction is retrieval-augmented generation (RAG)~\cite{lewis2020rag}: chunk the repository, embed every chunk, build a vector index, and retrieve the top-$k$ fragments most similar to the issue text. While effective, this pipeline carries costs that are seldom reported. Embedding an entire repository requires per-token API calls or dedicated GPU time; the resulting index must be stored and, in continuous-integration settings, rebuilt whenever the codebase changes. Most SWE-bench studies report \emph{resolve rate} as the sole metric, making it impossible to judge whether a method is practically deployable at scale. We argue that \textbf{context-construction cost deserves first-class status alongside resolve rate} when evaluating repository-level repair.

Programming languages, however, offer an alternative source of relevance signal that is both cheaper and more precise than semantic similarity: \emph{structural dependency}. Import statements, function calls, class inheritance, and module boundaries define an explicit graph that can be extracted with deterministic, non-LLM tooling such as Tree-sitter~\cite{brunsfeld2018treesitter} at negligible marginal cost. A dependency graph built once per repository snapshot can be traversed for every incoming issue without further embedding or indexing. Prior work on repository maps---notably Aider's \texttt{RepoMap}~\cite{aider2023}---has explored structural summaries for LLM context, but no controlled study has systematically compared graph-guided retrieval against RAG under identical LLM, patching, and evaluation conditions.

In this paper we introduce \emph{GraphManager} and its flagship strategy, \textbf{Graph-Map Progressive (GM-P)}. GM-P constructs a language-level dependency graph of the target repository, seeds the traversal from the suspected fault locale, and progressively expands context along structural edges---using agentic tool calls (\texttt{search\_nodes}, \texttt{get\_neighbors}, \texttt{get\_file\_summary})---until a token budget is filled. We evaluate GM-P against 10 alternative context-construction strategies spanning BM25 lexical search, fixed- and progressive-window RAG, agentic cold-start exploration, and controlled re-implementations inspired by Agentless-style localization~\cite{xia2024agentless} and Aider-style repository maps~\cite{aider2023}. All 11 methods share the same downstream LLM (Gemini~\cite{gemini2024}), the same patching pipeline, and the same evaluation harness across 500 SWE-bench Verified instances drawn from 12 open-source repositories. We report both resolve rate and per-resolved-issue token cost, and we assess pairwise differences with McNemar's test~\cite{mcnemar1947} corrected for multiplicity via the Holm--Bonferroni procedure.

Our principal findings are as follows. GM-P resolves 234 of 500 issues (46.8\%, 95\% CI [42.5\%, 51.2\%]), significantly outperforming every baseline in all nine pairwise comparisons after correction ($p < 0.05$). The oracle combination of all methods reaches 53.8\%, indicating that the primary ceiling is seed selection rather than retrieval depth. GM-P's cost of ${\sim}$256K tokens per resolved issue is approximately 7.5$\times$ lower than RAG-Progressive (1.93M tokens), making it the most cost-effective method among the high-performing tier.\footnote{Both the dependency graph and the RAG vector index were pre-built in our harness; the cost figures reflect method-specific LLM token consumption. We discuss this dual-build caveat in Section~\ref{sec:threats}.} Methods such as BM25 are cheaper in absolute terms but resolve substantially fewer issues. We emphasize that our baselines are \emph{controlled re-implementations} designed to isolate the effect of context-construction strategy; we do not claim to reproduce or surpass the full pipelines of Agentless, SWE-Agent, or Aider. The remainder of the paper is organized as follows: Section~\ref{sec:related} surveys related work; Section~\ref{sec:method} details the GraphManager framework and all compared strategies; Section~\ref{sec:setup} describes the experimental setup; Section~\ref{sec:results} presents results and statistical analysis; Section~\ref{sec:threats} discusses threats to validity; and Section~\ref{sec:conclusion} concludes.

%% ----------------------------------------------------------------
